\begin{definition}
    \emph{Поле} --- это коммутативное ассоциативное кольцо с единицей, в котором у
    каждого ненулевого элемента есть обратный.
\end{definition}

\section{Расширения полей}

\begin{definition}
    Пусть $F$ и $E$ --- поля, такие что $F$ является подполем $E$, то есть $F$ --- это подкольцо,
    которое само является полем. Тогда говорим, что имеем \emph{расширение полей} и пишем
    \[
        E / F.
    \]
\end{definition}

\begin{example}
    Примеры расширений:
    \begin{enumerate}
        \item $\mathbb{C} / \mathbb{R}$.
        \item $\mathbb{R} / \mathbb{Q}$.
        \item $\mathbb{Q}(\sqrt{2}) / \mathbb{Q}$, где $\mathbb{Q}(\sqrt{2}) = \mathbb{Q}[x] / (x^2 - 1)$.
    \end{enumerate}
\end{example}

\begin{stmt}
    Если $E / F$ --- расширение полей, то $E$ является векторным пространством над $F$.

    \begin{proof}
        $E$ замкнуто относительно сложения, ассоциативность и аддитивность выполнены.
        Элементы $F$ лежат в $E$ (как подполе), поэтому умножение на скаляры из $F$ определено
        как умножение в $E$. Все аксиомы векторного пространства выполнены.
    \end{proof}
\end{stmt}

\begin{definition}
    Расширение $E / F$ называется \emph{конечным}, если $E$ как векторное пространство над $F$
    имеет конечную размерность. Размерность обозначается
    \[
        [E : F] \coloneqq \dim_F E.
    \]
\end{definition}

\begin{example}
    Размерности расширений:
    \begin{enumerate}
        \item $[\mathbb{C} : \mathbb{R}] = 2$, базис $\{1, i\}$: любой $z \in \mathbb{C}$ записывается как $a \cdot 1 + b \cdot i$.
        \item $[\mathbb{R} : \mathbb{Q}] = \infty$: $\mathbb{Q}$ счётно, любая конечная линейная комбинация
        счётного числа элементов над $\mathbb{Q}$ даёт счётное множество, а $\mathbb{R}$ несчётно.
        \item $[\mathbb{Q}(\sqrt{2}) : \mathbb{Q}] = 2$, базис $\{1, \sqrt{2}\}$: любой элемент записывается как
        $a + b\sqrt{2}$, $a, b \in \mathbb{Q}$.
    \end{enumerate}
\end{example}

\begin{stmt}
    Пусть $K \subseteq F \subseteq E$ --- цепочка расширений полей. Тогда
    \[
        [E : K] < \infty \iff [F : K] < \infty \text{ и } [E : F] < \infty.
    \]
    В этом случае
    \[
        [E : K] = [E : F] \cdot [F : K].
    \]

    \begin{proof}
        Первая часть упражнение (следует из общих свойств векторных пространств).

        Для второй части построим базис $E$ над $K$ явно.
        Пусть $f_1, \ldots, f_n$ --- базис $F/K$, а $e_1, \ldots, e_m$ --- базис $E/F$.
        Утверждаем, что набор $\{f_i e_j\}_{i=1,\ldots,n;\, j=1,\ldots,m}$ (все попарные произведения) является базисом $E$ над $K$.

        \textbf{Порождаемость.} Возьмём $v \in E$. Так как $e_i$ --- базис $E$ над $F$:
        \[
            v = \sum_i \alpha_i e_i, \quad \alpha_i \in F.
        \]
        Так как $f_j$ --- базис $F$ над $K$:
        \[
            \alpha_i = \sum_j \beta_{ij} f_j, \quad \beta_{ij} \in K.
        \]
        Подставляя:
        \[
            v = \sum_i \alpha_i  e_i = \sum_i (\sum_j b_{ij} f_j)e_i = \sum_{i,j} \beta_{ij} (f_i e_j).
        \]

        \textbf{Линейная независимость.} Пусть $\sum_{i,j} \gamma_{ij} (e_i f_j) = 0$, $\gamma_{ij} \in K$.
        Сгруппируем по $e_i$:
        \[
            \sum_i \Bigl(\sum_j \gamma_{ij} f_j\Bigr) e_i = 0.
        \]
        Так как $(e_i)$ -- базис, то $\sum_j \gamma_{ij} f_j = 0 \quad \forall i \Rightarrow$ так как $(f_i)$ -- базис $F/K$,
        $\gamma_{ij} = 0 \quad \forall i,j$
    \end{proof}
\end{stmt}

\begin{remark}
    Конструктивно: базис большого расширения строится как набор всех
    попарных произведений базисов промежуточных расширений.
\end{remark}

\section{Подполе, порождённое множеством}

\begin{definition}
    Пусть $E / F$ --- расширение, $S \subseteq E$. \emph{Подполем, порождённым} $S$ над $F$,
    называется
    \[
        F(S) \coloneqq \bigcap_{\substack{T \text{ --- подполе } E \\ F \subseteq T,\; S \subseteq T}} T.
    \]
    Это наименьшее подполе $E$, содержащее $F$ и $S$.
\end{definition}

\begin{stmt}
    Пересечение двух подполей --- подполе.

    \begin{proof}
        Пусть $F_1, F_2$ --- подполя $E$. Возьмём $\alpha_1, \alpha_2 \in F_1 \cap F_2$. Так как они
        лежат в $F_1$, там лежат и $\alpha_1 + \alpha_2$, и $\alpha_1 \alpha_2$. Аналогично для $F_2$.
        Единица лежит в $F_1$ и $F_2$, значит и в пересечении. Обратный $\alpha_1^{-1}$ лежит в каждом
        (как поле), значит и в пересечении.
    \end{proof}
\end{stmt}

\begin{stmt}
    Пересечение любого (в том числе бесконечного) семейства подполей --- подполе.

    \begin{proof}
        Возьмём $\alpha_1, \alpha_2$ из пересечения. Они лежат в каждом $T_\lambda$, поэтому там
        лежат $\alpha_1 + \alpha_2$ и $\alpha_1 \alpha_2$ (каждое $T_\lambda$ --- подкольцо). Значит,
        они лежат и во всём пересечении. Обратный элемент --- аналогично.
    \end{proof}
\end{stmt}

\section{Конечно порождённые и простые расширения}

\begin{definition}
    Расширение $E / F$ называется \emph{конечно порождённым}, если существует конечное
    $S \subseteq E$ такое, что $E = F(S)$.
\end{definition}

\begin{remark}
    Конечная порождённость \emph{не эквивалентна} конечности расширения. Контрпример: $\mathbb{Q}(\pi) / \mathbb{Q}$.
    Оно конечно порождено (один порождающий $\pi$), но бесконечно как расширение:
    числа $1, \pi, \pi^2, \ldots$ линейно независимы над $\mathbb{Q}$ (так как $\pi$ трансцендентно),
    поэтому $[\mathbb{Q}(\pi) : \mathbb{Q}] = \infty$.
\end{remark}

\begin{definition}
    Пусть $E / F$ --- расширение, $\alpha \in E$. Элемент $\alpha$ называется \emph{алгебраическим
    над $F$}, если существует ненулевой многочлен $p(x) \in F[x]$ такой, что $p(\alpha) = 0$.
    Иначе $\alpha$ называется \emph{трансцендентным над $F$}.

\end{definition}

\begin{definition}
    Расширение $E / F$ называется \emph{алгебраическим}, если все его элементы алгебраичны над $F$.
\end{definition}

\begin{examplelist}
    \item $\mathbb{C} / \mathbb{R}$ алгебраично: для любого $z \in \mathbb{C}$
    многочлен $(x - z)(x - \bar{z}) = x^2 - 2\mathrm{Re}(z)\cdot x + |z|^2$ имеет вещественные коэффициенты и корень $z$.
    \item $\mathbb{Q}(\sqrt{2}) / \mathbb{Q}$ алгебраично: для $z = a + b\sqrt{2}$
    многочлен $(x - z)(x - \bar{z})$ с $\bar{z} = a - b\sqrt{2}$ вещественный, а именно
    $x^2 - 2ax + (a^2 - 2b^2) \in \mathbb{Q}[x]$.
\end{examplelist}


\begin{definition}
    Расширение $E / F$ называется \emph{простым}, если $E = F(\alpha), \alpha \in E$.
\end{definition}

\begin{examplelist}
    \item $\mathbb{C} / \mathbb{R}$: $\mathbb{C} = \mathbb{R}(i)$.
    \item $\mathbb{F}_{25} / \mathbb{F}_5$.
    \item $\mathbb{Q}(\sqrt{2}, \sqrt{3}) / \mathbb{Q}: \mathbb{Q}(\sqrt{2}, \sqrt{3}) = \mathbb{Q}(\sqrt{2} + \sqrt{3})$.
    \begin{proof}
        $\newline \supseteq$: очевидно \\
        $\subseteq$:
        \begin{equation*}
            (\sqrt{2} + \sqrt{3})^{-1} = \frac{1}{\sqrt{2} + \sqrt{3}} = \frac{\sqrt{2} - \sqrt{3}}{-1} = \sqrt{3} - \sqrt{2}
        \end{equation*}

        \begin{equation*}
            \frac{(\sqrt{3} + \sqrt{2}) + (\sqrt{3} - \sqrt{2})}{2} = \sqrt{3} \quad \quad (\sqrt{3} + \sqrt{2}) - \sqrt{3} = \sqrt{2}
        \end{equation*}
    \end{proof}
\end{examplelist}



\begin{stmt}


Пусть $E/F$ -- расширение, $\alpha \in E$.
\[
    \varphi\colon F[x] \to \underbrace{F[\alpha]}_{=\{\sum \gamma_i \alpha^i | \gamma_i \in F\}} \qquad p \mapsto p(\alpha).
\]
Это гомоморфизм колец (сюръекция на $F[\alpha]$ --- подкольцо $E$, порождённое степенями $\alpha$).
Возможны два случая.

\textbf{Случай 1: $\ker\varphi = \{0\}$.} Ни один ненулевой многочлен не зануляется на $\alpha$.
По определению, $\alpha$ трансцендентно.

\textbf{Случай 2: $\ker\varphi \neq \{0\}$.} Ядро --- ненулевой идеал в $F[x]$.
Так как $F[x]$ --- область главных идеалов, $\ker\varphi = (f)$ для некоторого многочлена $f(x)$.

\begin{definition}
    Многочлен унитарный, если его старший коэффициент равен $1$.
\end{definition}

\begin{remark}
    Всегда считаем $f$ унитарным.
\end{remark}

\[
\varphi(f) = f(\alpha) = 0 \Leftrightarrow \alpha \text{ --- алгебраический}.
\]
\end{stmt}

\begin{definition}
    Многочлен $f \in F[x]$, порождающий $\ker\varphi$, называется \emph{минимальным многочленом}
    для $\alpha$ над $F$.
\end{definition}

\begin{remark}
    $\varphi \colon F[x] \to F[\alpha]$ --- сюръекция. \\
    Пусть у нас случай 2. По теореме о гомоморфизме:
    
    \begin{align*}
        \frac{F[x]}{\ker \varphi} &= \frac{F[x]}{(f(x))} \cong F[\alpha]
    \end{align*}

    \[
    \frac{F[x]}{(f(x))} \text{ поле. А значит $F[\alpha]$ --- тоже поле.}
    \]
    \begin{cor}
        Если $\alpha$ алгебраично над $F$, то при записи простого расширения не нужно насильно
        добавлять обратные элементы --- они добавляются автоматически:
        \[
            F[\alpha] = F(\alpha).
        \]
        В частности: $\mathbb{Q}(\sqrt{2}) = \mathbb{Q}[\sqrt{2}]$. \\
        Для трансцендентных элементов это неверно: $F[\pi] \not\cong F(\pi)$ (первое --- кольцо, не поле).
    \end{cor}
\end{remark}







\section{Конечность, алгебраичность и конечная порождённость}

 \begin{theorem}
    Расширение $E / F$ конечно тогда и только тогда, когда оно конечно порождено и алгебраично.

\begin{proof}
    \textbf{($\Rightarrow$)} Пусть $[E : F] = n < \infty$. 
    
    1. \emph{Конечная порождённость:} Возьмём базис $e_1, \ldots, e_n$ пространства $E$ над $F$. 
    Тогда любой элемент из $E$ представим в виде $\sum \gamma_i e_i$, значит $E = F(e_1, \ldots, e_n)$.
    Тот факт, что элементы порождают $E$ как векторное пространство, автоматически означает, что они порождают его и как расширение поля.
    
    2. \emph{Алгебраичность:} Пусть $t \in E$. Рассмотрим систему элементов $1, t, t^2, \ldots, t^n$.
    Это $n+1$ вектор в $n$-мерном пространстве, следовательно, они линейно зависимы над $F$:
    \[
        \exists\, \gamma_0, \gamma_1, \ldots, \gamma_n \in F \text{ (не все равные 0), такие что } \sum_{i=0}^n \gamma_i t^i = 0.
    \]
    Это означает, что $t$ является корнем ненулевого многочлена из $F[x]$, то есть $t$ алгебраичен над $F$.

    \textbf{($\Leftarrow$)} Пусть $E = F(t_1, \ldots, t_m)$, где каждый $t_i$ алгебраичен над $F$.
    Докажем индукцией по числу порождающих $m$, что $[E : F] < \infty$.

    \textbf{База $m = 1$:} $E = F(t_1)$. Так как $t_1$ алгебраичен над $F$,
    то степень расширения равна степени его минимального многочлена: $[F(t_1) : F] < \infty$.

    \textbf{Переход:} Пусть для $m$ порождающих утверждение верно. Рассмотрим $E = F(t_1, \ldots, t_m, t_{m+1})$.
    Обозначим промежуточное поле $L = F(t_1, \ldots, t_m)$. 
    По предположению индукции $[L : F] < \infty$. 
    Элемент $t_{m+1}$ алгебраичен над $F$, следовательно, он тем более алгебраичен над $L$
    (его минимальный многочлен над $L$ является делителем минимального многочлена над $F$). 
    
    Тогда имеем цепочку расширений:
    \[
        F \xhookrightarrow{\text{кон}} L \xhookrightarrow{\text{кон}} L(t_{m+1}) = E
    \]
    Так как $E = L(t_{m+1})$ — простое алгебраическое расширение поля $L$, то $[E : L] < \infty$. По теореме о башне расширений:
    \[
        [E : F] = [E : L] \cdot [L : F]
    \]
    Произведение двух конечных чисел конечно, значит $[E : F] < \infty$.
\end{proof} 
\end{theorem}

\begin{stmt}
    Пусть $E/F$ --- расширение полей, и в $F[t_1, \dots, t_n] \subseteq E$ все элементы $t_i$ алгебраичны над $F$. Тогда $F[t_1, \dots, t_n] = F(t_1, \dots, t_n)$.

    \begin{proof}
        Индукция по $n$.
            \textbf{База:} $n=1$. Имеем простое расширение $F[t_1]/F$. Известно, что для алгебраического элемента $F[t_1] = F(t_1)$.
            \textbf{Переход:}
            \[
            F[t_1, \dots, t_{n+1}] = F[t_1, \dots, t_n][t_{n+1}] = F(t_1, \dots, t_n)[t_{n+1}] = F(t_1, \dots, t_{n+1})
            \]
    \end{proof}
\end{stmt}


\begin{stmt}
    Алгебраичность транзитивна: если $F \subseteq L \subseteq E$, расширение $L / F$ алгебраическое
    и $E / L$ алгебраическое, то $E / F$ алгебраическое.

\begin{proof}
    Возьмём произвольный $\alpha \in E$. Так как $E / L$ алгебраическое, существует ненулевой многочлен, зануляющий $\alpha$:
    \[
        f(x) = x^n + \beta_1 x^{n-1} + \dots + \beta_n \in L[x], \quad f(\alpha) = 0.
    \]
    Все коэффициенты $\beta_1, \dots, \beta_n$ принадлежат $L$, а значит, алгебраичны над $F$ (так как $L / F$ алгебраическое).
    Рассмотрим поле $K = F(\beta_1, \dots, \beta_n)$. 
    
    1. Расширение $K/F$ конечно, так как оно конечно порождено алгебраическими элементами $\beta_i$.
    2. Элемент $\alpha$ алгебраичен над $K$, так как $f(x) \in K[x]$. Следовательно, простое расширение $K(\alpha)/K$ конечно.

    Рассмотрим цепочку конечных расширений:
    \[
        F \xhookrightarrow{\text{кон}} F(\beta_1, \dots, \beta_n) \xhookrightarrow{\text{кон}} F(\beta_1, \dots, \beta_n, \alpha).
    \]
    По теореме о башне расширений степень итогового расширения есть произведение степеней:
    \[
        [F(\beta_1, \dots, \beta_n, \alpha) : F] = [F(\beta_1, \dots, \beta_n, \alpha) : K] \cdot [K : F] < \infty.
    \]
    Так как расширение $F(\beta_1, \dots, \beta_n, \alpha) / F$ конечно, оно является алгебраическим. 
    Следовательно, элемент $\alpha$ алгебраичен над $F$. В силу произвольности выбора $\alpha$, всё расширение $E/F$ алгебраично. \qed
\end{proof}
\end{stmt}


\section{Поле разложения}

\begin{definition}
    Пусть $F$ --- поле, $p \in F[x]$ --- многочлен. \emph{Полем разложения} $p(x)$ называется
    поле $E \supseteq F$ такое, что $p$ раскладывается на произведение линейных множителей в $E[x]$:
    \[
        p(x) = c \prod_{i=1}^n (x - \alpha_i), \quad \alpha_i \in E.
    \]
\end{definition}

\begin{examplelist}
        \item $p(x) = x^2 + 1 \in \mathbb{R}[x]$. Корни $\pm i$, поле разложения: $\mathbb{C} = \mathbb{R}(i)$.
        \item $p(x) = x^2 + 1 \in \mathbb{Q}[x]$. Поле разложения: $\mathbb{Q}(i)$.
        \item $p(x) = x^2 - 2 \in \mathbb{Q}[x]$. Поле разложения: $\mathbb{Q}(\sqrt{2})$.
\end{examplelist}

\section{Алгебраически замкнутые поля}

\begin{theorem}
    Для поля $E$ следующие условия равносильны:
    \begin{enumerate}
        \item Любой неконстантный $p \in E[x]$ раскладывается над $E$ в произведение линейных множителей.
        \item Любой неконстантный $p \in E[x]$ имеет хотя бы один корень в $E$.
        \item Все неприводимые многочлены из $E[x]$ имеют степень $1$.
        \item $\forall$ поле $T \supset E$ и которое конечно надо $E$, рано $E$.
    \end{enumerate}

\begin{proof}
        $(1 \Rightarrow 2)$: Ясно: линейные множители вида $(x - \alpha)$ сразу дают корни.

        $(2 \Rightarrow 3)$: Пусть $p \in E[x]$ неприводим.
        Если $\deg p \geq 2$, то по (2) у $p$ есть корень $\gamma$,
        следовательно $p \vdots (x - \gamma)$, что противоречит неприводимости.

        $(3 \Rightarrow 2)$: Пусть $p(x)$ не имеет корней.
        Возьмём его неприводимый множитель $q(x)$, у которого тоже нет корней.
        Но по (3) $\deg q = 1$, а любой многочлен первой степени имеет корень --- противоречие.

        $(2 \Rightarrow 1)$: Если $p(\alpha) = 0$, то $p(x) \vdots (x - \alpha)$.
        Применяем то же самое рассуждение к частному $p(x)/(x - \alpha)$ и продолжаем до полного разложения.

        $(1 \Rightarrow 4)$: Пусть $T / E$ --- конечное расширение и $\alpha \in T$.
        Тогда $\alpha$ алгебраичен над $E$, то есть $\exists f(x) \in E[x]: f(\alpha) = 0$.
        По (1) $f(x) = c(x - \gamma_1)\dots(x - \gamma_n)$, где все $\gamma_i \in E$.
        Значит $\alpha = \gamma_i$ для некоторого $i$, откуда $\alpha \in E$. Следовательно, $T = E$.

        $(4 \Rightarrow 2)$: Возьмём неприводимый $f(x) \in E[x]$.
        Рассмотрим поле $E[x]/(f(x)) \supset E$, которое является конечным расширением.
        По (4) это поле совпадает с $E$. Рассмотрим каноническую проекцию:
        \[
            \pi: E[x] \to E[x]/(f(x)), \quad x \mapsto \pi(x)
        \]
        Так как $\pi(f(x)) = 0$, получаем $f(\pi(x)) = 0$.
        При этом $\pi(x) \in E$, значит у $f$ есть корень в $E$.
    \end{proof}
\end{theorem}

\begin{definition}
    Поле $E$, удовлетворяющее одному (и значит всем) условиям теоремы, называется
    \emph{алгебраически замкнутым}.
\end{definition}
